\documentclass[10pt, a4paper]{article}

% Parametri che modificano il file main.tex
% Le uniche parti da cambiare su main.tex sono:
% - tabella versioni
% - testo file

\def\titolo{Analisi dei \\ \ requisiti} % \\ per andare a capo, \ per spazio
\def\titoloHeader{Analisi dei requisiti}
\def\data{}

% \def\listaComponenti{
% Bresolin G.,
% Campese M.,
% Ciriolo I.,
% Dugo A.,
% Feltrin E.,
% Michelon R.,
% Orlandi G.
% }

\usepackage{style}
\usepackage{headerfooter}

\title{\data}
\author{SWEetCode}

\begin{document}

% PRIMA PAGINA
\begin{titlepage}
    \thispagestyle{empty}
    \begin{tikzpicture}[remember picture, overlay]
        % TRIANGOLI
        \draw[fill=secondarycolor, secondarycolor] (current page.north west) -- (current page.south west) -- (8.8, -28);
        \draw[fill=primarycolor, primarycolor] (-3, 5) -- (4, -13.6) -- (11, 5);

        % LOGO
        \node [xshift=-5cm, yshift=25cm] (logo) at (current page.south east) {\includegraphics[width=6.5cm]{img/logo.png}};

        % SWEETCODE - DATE
        \node [anchor=north east, align=right, xshift=-1.2cm, yshift=20.5cm, text=black] (sweetcode) at (current page.south east) {\fontsize{32pt}{36pt}\selectfont SWEetCode};
        \draw[line width=4pt, lightcol] ([xshift=-3cm, yshift=-0.37cm]sweetcode.south west) -- ([yshift=-0.37cm]sweetcode.south east);
        \node [anchor=north east, align=right, xshift=-1.2cm, yshift=18.7cm, text=black] (date) at (current page.south east){\fontsize{24pt}{24pt} \selectfont \data};

        % NOME FILE
        \node [anchor=north east, text width=15cm, align=right, xshift=-1.2cm, yshift=17cm, text=black] (titolo) at (current page.south east){\fontsize{48pt}{48pt}\textbf{\titolo}};

        % BOX DATI PARTECIPANTI
        \node[anchor=north east, xshift=-1.2cm, yshift=12cm, minimum width=8cm] (box) at (current page.south east){};

        % COMPONENTI DEL GRUPPO
        \node[anchor=north west, yshift=-1cm, align=left] (dati1) at (box.north west) {\fontsize{16pt}{16pt}\selectfont \textbf{Componenti del gruppo}};
        \draw[line width=4pt, lightcol] (dati1.south west) -- ([xshift=8cm]dati1.south west);
        
        \node[anchor=north west, align=left] (dati11) at (dati1.south west)
        {\fontsize{14pt}{14pt}\selectfont Bresolin G.};
        \node[anchor=north west, yshift=-0.7cm, align=left] (dati12) at (dati1.south west)
        {\fontsize{14pt}{14pt}\selectfont Campese M.};
        \node[anchor=north west, yshift=-1.4cm, align=left] (dati13) at (dati1.south west)
        {\fontsize{14pt}{14pt}\selectfont Ciriolo I.};
        \node[anchor=north west, yshift=-2.1cm, align=left] (dati14) at (dati1.south west)
        {\fontsize{14pt}{14pt}\selectfont Dugo A.};
        \node[anchor=north west, yshift=-2.8cm, align=left] (dati15) at (dati1.south west)
        {\fontsize{14pt}{14pt}\selectfont Feltrin E.};
        \node[anchor=north west, yshift=-3.5cm, align=left] (dati16) at (dati1.south west)
        {\fontsize{14pt}{14pt}\selectfont Michelon R.};
        \node[anchor=north west, yshift=-4.2cm, align=left] (dati17) at (dati1.south west)
        {\fontsize{14pt}{14pt}\selectfont Orlandi G.};
        
        % UNIPD - SWE
        \node [xshift=4.4cm, yshift=2.3cm, draw, secondarycolor, text=white] (uni) at (current page.south west) {\fontsize{20pt}{20pt} \selectfont Università di Padova};
        \node [xshift=0.65cm, yshift=0.7cm, draw, secondarycolor, text=white, below=of uni] (corso) {\fontsize{20pt}{20pt}\selectfont Ingegneria del Software};

        % FIRMA
        % \draw[line width=4pt, lightcol] ([xshift=-1.2cm, yshift=1.8cm]current page.south east) -- ([xshift=-8cm, yshift=1.8cm]current page.south east);
        % \node[anchor=north west, xshift=12.9cm, yshift=1.45cm, align=left] at (current page.south west)
        % {\fontsize{13pt}{13pt}\selectfont L'Amministratore: Feltrin E.};
        
    \end{tikzpicture}
\end{titlepage}

% INIZIO PAGINE
\columnratio{0.31}
\setlength{\columnsep}{2.2em}
\setlength{\columnseprule}{4pt}
\colseprulecolor{lightcol}
\begin{paracol}{2}

% 11111111111111111111111111111111111111111111111111111111111111111111111111111

\intestazione
\vspace{5.0em}

\partecipanti

\newpage

\switchcolumn
\revisioneAzioni
\vspace{23.5em}

\ordineGiorno

\newpage

% 22222222222222222222222222222222222222222222222222222222222222222222222222222

\switchcolumn

\heading{Discussione}\\
\textbf{Sintesi degli argomenti\\discussi}

~\newpage
~ \\\\\\\\\\\\\\\\\\\\\\\\\\\\\\\\\\\\\\~

\heading{Decisioni}\\
\textbf{Decisioni prese durante\\la discussione}

\switchcolumn

\discussione
Successivamente alla discussione, i rappresentanti dell'azienda hanno chiarito alcuni aspetti necessari all'avviamento del progetto. Viene avviato un canale \textit{Slack} per le comunicazioni con l'azienda, che mette inoltre a disposizione un \textit{repository} \textit{GitHub}. In aggiunta viene messo a disposizione un calendario \textit{Google} condiviso per la programmazione delle riunioni.

L'azienda spiega che il primo blocco di lavoro sarà di 6 settimane, entro le quali sarà richiesto un PoC (con scadenza entro il 2023-12-06), nel quale deve essere presente una solida analisi dei requisiti. Il secondo blocco di lavoro (MVP) sarà invece di circa 14 settimane, con task e scadenze ancora da definire.

Viene infine specificato che verrà richiesto di elaborare un'analisi dei costi di costruzione e di mantenimento del progetto.

\vspace{12.5em}

\decisioni

\end{paracol}

\newpage

\heading{Azioni da Intraprendere}

{\renewcommand{\arraystretch}{1.5}
\begin{tabularx}{\textwidth}{c|X|c|c|c}
\textbf{Task Jira} & \textbf{Azione} & \textbf{Incaricato} & \textbf{Revisore} & \textbf{Scadenza} \\
\hline
\#35 &
Stesura verbale riunione &
Ciriolo I. &
Campese M. &
$2023-11-13$ \\
\hline
- &
Stesura prima bozza di Analisi dei requisiti &
Ciriolo I. &
Campese M. &
$2023-11-15$ \\
\hline
- &
Discussione riguardo l'organizzazione pubblica/privata del \textit{repository} &
Team &
Team &
$2023-11-15$ \\
\hline
- &
Iscrizione al canale \textit{Slack} ed al \textit{repository} forniti dall'azienda &
Feltrin E. &
Campese M. &
$2023-11-15$ \\
\hline
- &
Invio all'azienda di account e-mail personali dei membri del team &
Feltrin E. &
Campese M. &
$2023-11-15$ \\
\end{tabularx}}

\vspace{3em}

\heading{Altro}

\textbf{Prossima Riunione}

La prossima riunione con AzzurroDigitale è pianificata il 2023-11-15 alle 15:00 per la discussione dell'Analisi dei requisiti.

\end{document}